\chapter{Implementation}\label{chap:implementation}

\section{Introduction}

This chapter explains how the design of \cref{chap:design} was realised. It describes the technology stack and the
structure of the repository, then follows the implementation through the backend, the mobile app, the web console
and the AI components, and finishes with the integrations, deployment and the main engineering challenges. Short
excerpts of the real source code illustrate the points that matter most; they are taken unchanged from the
repository, with only blank lines and long comments removed.

\section{Technology Stack}

The stack (\cref{tab:stack}) uses mainstream, well-documented technologies with strong ecosystems, chosen so that
the team could deliver a complete system and so that the project can be maintained afterwards.

\begin{table}[!htbp]
\caption{Technology stack.}\label{tab:stack}
\centering\tabfont\rowband
\begin{tabular}{@{}>{\bfseries\raggedright\arraybackslash}p{2.6cm}>{\raggedright\arraybackslash}p{12cm}@{}}
\rowcolor{najdanavy}\textcolor{white}{\sffamily\bfseries\small Layer} & \textcolor{white}{\sffamily\bfseries\small Technology} \\
Backend & Java 21; Spring Boot 4.1 (Web, Security, Data JPA, Validation, WebSocket/STOMP) \cite{springboot}; Hibernate; Flyway \cite{flyway}; Firebase Admin SDK; PostgreSQL 17 \cite{postgresql} \\
Mobile app & Flutter \cite{flutter} with Riverpod \cite{riverpod} and go\_router; Dio; stomp\_dart\_client; firebase\_auth, firebase\_messaging, google\_sign\_in; geolocator; maplibre\_gl and flutter\_map; image\_picker, record, audioplayers, video\_player; local\_auth; connectivity\_plus; shared\_preferences \\
Web console & Next.js 16 with the App Router and React 19 \cite{nextjs}; TanStack Query 5 \cite{tanstackquery}; Zustand 5; next-intl 4; Firebase web SDK; @stomp/stompjs; MapLibre GL with react-map-gl; idb-keyval; EmailJS \\
AI model service & Python 3.12; FastAPI \cite{fastapi}; Uvicorn; pandas and NumPy; scikit-learn \cite{pedregosa2011sklearn}; XGBoost \cite{chen2016xgboost}; joblib \\
External services & Firebase Authentication and Cloud Messaging \cite{firebase}; Supabase Storage \cite{supabase}; OpenRouter \cite{openrouter}; Brevo \cite{brevo}; Didit \cite{didit}; Nominatim \cite{nominatim}; Overpass; OSRM \cite{osrm}; OpenFreeMap; GitHub Releases \\
Delivery & Docker and Docker Compose \cite{docker}; GitHub Actions \cite{githubactions}; Vercel (web); container hosting (API); signed Android APK \\
\end{tabular}
\end{table}

\section{Repository Structure}

The project is a single repository with one folder per deployable unit:

\begin{lstlisting}[language={},numbers=none,caption={Repository layout.},label={lst:repo}]
najda/
|-- backend/              Spring Boot API: modules, Flyway migrations, tests
|-- mobile/               Flutter app (Android): features, core services, tests
|-- web/                  Next.js console: app routes, components, hooks, store
|-- ai/                   FastAPI model service: registry, tests, Dockerfile
|-- docs/                 documentation, diagrams and their UML sources
|-- docker-compose.yml    local stack: PostgreSQL, AI service, backend, web
`-- .github/workflows/    CI for the AI service and the mobile app, release build
\end{lstlisting}

\section{Backend Implementation}

\subsection{Layering and conventions}
Each feature module has the same three layers. \emph{Controllers} translate HTTP to method calls and carry the
role annotations (\code{@PreAuthorize}); \emph{services} contain the rules and own the transactions; \emph{repositories}
are Spring Data interfaces. Request and response bodies are immutable record types, so an entity is never exposed
directly. Rule violations are raised as typed exceptions (not found, conflict, forbidden, too many requests, KYC or
legal-consent required) which a central handler maps to the right HTTP status and a stable error code that the clients
can act on.

\subsection{Concurrency control}\label{sec:impl-concurrency}
Two users can easily try to act on the same thing at the same time: two dispatchers offering the same unit, or a
dispatcher cancelling a mission as the crew accepts it. Unit assignment takes a \emph{pessimistic row lock} on the unit
(\cref{lst:assign}), so the availability check and the status change are one atomic step and the loser of a race
receives a clear conflict message. Incidents and missions carry a \code{@Version} field, so concurrent edits fail instead
of silently overwriting each other.

\lstinputlisting[language=Java,caption={Unit assignment: lock, guard, then create the mission (\code{MissionServiceImpl}, abridged).},label={lst:assign}]{code/MissionAssign.java}

Note the guard against a subtle failure of the domain: a responder can report an emergency too, and the unit they
crew must never be sent to their own report.

\subsection{Idempotent submission}
The client generates a request id for each report and sends it with every attempt. The service looks the id up
\emph{first}; if the incident exists it is returned unchanged, with no new side effects. A partial unique index makes the
guarantee hold even if two attempts race (\cref{lst:idempotency}).

\lstinputlisting[language=SQL,caption={The database-level guarantee behind idempotent submission.},label={lst:idempotency}]{code/idempotency.sql}

\subsection{Side effects after commit}\label{sec:impl-aftercommit}
Live updates and notifications are published from inside business transactions, but must reach clients only after the
data is committed; otherwise a client that refetches on a nudge reads the old state, or is told about an action that
rolled back. A tiny utility registers the action with Spring's transaction synchronisation (\cref{lst:aftercommit}) and
every event publisher goes through it.

\lstinputlisting[language=Java,caption={Running an action after the surrounding transaction commits (\code{AfterCommit}).},label={lst:aftercommit}]{code/AfterCommit.java}

\subsection{Caller trust as a pure function}
The caller-trust policy is implemented as a final class of static functions with no database access and no clock, so it
is trivial to test exhaustively (\cref{lst:trust}). Weights, half-life, window, tier thresholds and the cool-down
gap are named constants that mirror \cref{tab:trust-params}.

\lstinputlisting[language=Java,caption={Scoring, decay and tiers in \code{CallerTrustPolicy}.},label={lst:trust}]{code/CallerTrustPolicy.java}

\subsection{One-time codes}
Registration codes are six random digits from a cryptographically secure generator. Only a keyed hash bound to the
e-mail address is stored, so a leaked row is useless without the server's secret, and verification compares hashes in
constant time (\cref{lst:otp}). A code is valid for 10 minutes, allows five attempts and can be re-sent only after 60 seconds;
code requests and verifications are also rate-limited per client address.

\lstinputlisting[language=Java,caption={Keyed hashing of one-time codes (\code{OtpCodec}).},label={lst:otp}]{code/OtpCodec.java}

\subsection{Persistence}
The schema is created by a Flyway baseline migration and validated against the entities by Hibernate at start-up, so
drift between code and database fails fast. Sensitive fields --- chat content, text evidence, free-text medical
fields, emergency-contact names and phones, and vitals notes --- are stored through attribute converters that encrypt
them with a configured key. Notifications and device tokens older than 90 days are purged nightly, and other
maintenance jobs run on fixed schedules (\cref{tab:jobs}).

\subsection{Evidence storage}
Evidence and application documents are stored in a private bucket. For an upload the service checks that the caller
owns a live incident and that the file extension is on an allow-list, then returns a short-lived signed upload URL;
the client uploads directly to storage and then registers the file, and the service confirms that the object exists
before saving the reference. Downloads use short-lived signed URLs issued only to people allowed to see the incident.
Deleting an incident removes its storage folder after the transaction commits.

\section{Mobile Application}

The Flutter app is organised into feature folders --- \code{auth}, \code{account}, \code{incident}, \code{map},
\code{chat}, \code{responder}, \code{handoff}, \code{kyc}, \code{legal}, \code{medical}, \code{notifications},
\code{guidance}, \code{applications}, \code{settings} and \code{status} --- over a \code{core} layer for networking, live
updates, push, storage, localisation, routing, security and theming. State is held in Riverpod providers; navigation uses
go\_router with guards that route a user by role, by profile completeness, by legal consent and by the app lock.

\subsection{Offline-first reporting}
The outbox is the most safety-critical piece of the client. When a report cannot be sent, it is first written to disk
together with its media, so a crash or a dead battery cannot lose it (\cref{lst:outbox}). Every attempt carries the same
request id, so a retry after a lost response returns the original incident. Progress is saved after creating the incident
and after each upload, so a retry resumes where it stopped; failed attempts back off from 30 seconds to 5 minutes, and
no attempt is made while the device is offline. Reports saved long ago ask the citizen to confirm before they are sent, and
reports the server refuses are shown to the citizen instead of being retried silently.

\lstinputlisting[language=Dart,caption={The guarantees of the offline outbox (\code{OutboxController}).},label={lst:outbox}]{code/outbox_controller.dart}

\subsection{Live updates, push and security}
A single STOMP connection delivers live changes; screens refetch when they are told something changed. Firebase Cloud
Messaging delivers push notifications on a dedicated channel for urgent types, with the token registered together with the
device language so that text arrives in the right language. An optional biometric or device-credential lock protects the
app when it returns from the background, and an update notifier offers newer releases found by the API on GitHub.

\subsection{Maps and location}
Maps are drawn with MapLibre using OpenFreeMap tiles, and routes and ETAs come from OSRM. A boundary check in the app
mirrors the server's Egypt boundary so that a citizen is told immediately when a location cannot be used.

\section{Web Console}

The console is a Next.js application with one route tree per workspace: dispatch, responder, hospital, facility manager and
administration, plus the citizen emergency, account, identity-verification and notification pages and the public
pages (landing, terms, privacy, status and support). Server state is cached by TanStack Query and invalidated by STOMP
subscriptions, which gives a live console without polling. The dispatcher workspace ranks available units by distance
and estimated travel time computed in the browser from live positions, and draws routes with OSRM; the hospital workspace
shows incoming transfers on a map with ETA and a live vitals stream. Text is localised in English and Arabic with
right-to-left layout through next-intl, and the interface supports light and dark themes.

\Needspace{16\baselineskip}
\section{AI Implementation}

\subsection{Priority and duplicate classification}
The priority service builds a structured prompt from the category, the optional message and the context, calls the
language model with temperature zero and a time-out, and parses the answer into a priority and a confidence. The
duplicate service first selects candidate incidents with a database query (same category, open, within 2\,km and 120
minutes) and asks the model only to judge those. Both run asynchronously after the transaction commits and both swallow
and log failures, so a provider outage never affects a report.

\subsection{Model-serving platform}
The model service loads each model package in two steps that guard against silent errors: it verifies the artifact's hash
against the manifest and refuses to load a model whose key libraries differ from the training environment
(\cref{lst:registry}). Prediction validates the input against the manifest --- unknown features, out-of-range values and
unknown categories are rejected --- and returns a label, a confidence and the full probability distribution.

\lstinputlisting[language=Python,caption={Loading a model package safely (\code{registry.py}).},label={lst:registry}]{code/registry_load.py}

\section{Integrations}

\paragraph{Firebase.} The backend verifies ID tokens with the Admin SDK, creates accounts (employees with a temporary
password and a one-time reset link), writes the role into a custom claim, reads the verified phone and e-mail during
synchronisation, and sends push notifications.
\paragraph{Didit.} A hosted verification session is created for the user; the signed webhook is authenticated by signature
and timestamp, de-duplicated by event id, and confirmed against the provider's decision before the status changes. A stale
status is reconciled directly with the provider when the app asks for it.
\paragraph{Supabase Storage.} Signed upload and download URLs for a private bucket; the service never handles file bytes.
\paragraph{Brevo.} Sends registration codes; its failure only fails the code request.
\paragraph{OpenStreetMap services.} Nominatim for reverse geocoding (throttled), Overpass (with mirror fall-back) for importing
hospitals, fire and police stations, and OSRM for routes. A bundled list of about 800 Egyptian facilities is loaded on first
start so the system is usable immediately.

\section{Deployment and Delivery}

The backend and the AI service are packaged as container images; the backend image runs on a Java 21 runtime as a
non-root user, and the AI service publishes no port. A single Compose file starts PostgreSQL, the AI service, the backend
and the web console in the right order, waiting for the database to be healthy. GitHub Actions runs the AI service tests
and image build on every push, analyses the mobile app on every pull request, and on a version tag builds and publishes the
signed production APK (\cref{tab:pipelines}). Configuration is entirely environment-driven (Appendix~\ref{app:configuration}).

\section{Engineering Challenges and Solutions}

\Cref{tab:challenges} lists the main engineering problems met during the implementation and how each was solved.

\begin{table}[!htbp]
\caption{Main engineering challenges and their solutions.}\label{tab:challenges}
\centering\tabfont\rowband
\begin{tabular}{@{}>{\raggedright\arraybackslash}p{4.6cm}>{\raggedright\arraybackslash}p{9.6cm}@{}}
\rowcolor{najdanavy}\textcolor{white}{\sffamily\bfseries\small Challenge} & \textcolor{white}{\sffamily\bfseries\small Solution} \\
An account lives in two systems (identity provider and database) that cannot share a transaction. & Verify the e-mail first; create the identity account; save the database rows in one transaction; delete the identity account if the database step fails. \\
A flaky network sends the same report twice. & Client-generated request id, checked before any other rule and backed by a partial unique index; the outbox reuses the same id. \\
Two dispatchers offer the same unit. & Pessimistic row lock on the unit and a clear conflict message. \\
Live updates arrive before the data is committed. & All broadcasts and notifications go through \code{AfterCommit}. \\
Webhooks arrive late, twice or forged. & Signature and timestamp check, event-id table, decision re-fetched from the provider, and a settled status is never undone by a late \emph{in progress} message. \\
Medical data is sensitive but life-saving. & Layered rule: report consent, standing consent, open incident, treating role; administrators excluded; every read audited and failing closed. \\
Abuse must be limited without ever refusing a caller. & A pure trust policy that only adds a short gap between reports and always shows the emergency numbers. \\
The language model may fail or be slow. & Asynchronous, time-limited, deterministic calls whose failure is logged and ignored; retry available. \\
Arabic and English with right-to-left layout everywhere, including notifications. & Localisation in both clients; notification text produced in both languages and delivered in the device's language. \\
\end{tabular}
\end{table}

\section{The User Interface}

The following screens illustrate the main workflows; further screens are in Appendix~\ref{app:ui}.

\screenshotrow[0.78\textheight]{mobile_report,mobile_tracking_and_chat,mobile_handoff,mobile_handoff2}{Mobile app, full-length screens from left to right: the emergency report screen with the first-aid guidance below it; following the response (status timeline, map, assigned responder, incident details and chat); the ambulance handoff on arrival (patient information, medical profile and choice of destination hospital); and the handoff while transporting (vital-signs entry and the latest vitals sent).}{fig:ui-mobile}
\screenshot[0.92\linewidth]{web_dispatch}{Web console: the dispatcher workspace with live queue, map and unit panel.}{fig:ui-dispatch}

\section{Chapter Summary}

The implementation realises each design decision in a small, testable piece of code: a row lock for assignment, a unique index
for idempotency, an after-commit helper for events, a pure function for caller trust, keyed hashes for codes, and a
durable outbox on the phone. Together with the integrations and the deployment pipeline they form a complete platform;
\cref{chap:testing} describes how it is verified.
